\newcommand{\DoNotLoadEpstopdf}{}
\pdfinfoomitdate=1
\pdftrailerid{}
\pdfsuppressptexinfo=-1
\documentclass[11pt,letterpaper]{article}
\usepackage[T1]{fontenc}
\usepackage{lmodern}
\usepackage{amsmath,amssymb,amsthm,mathtools}
\usepackage[margin=1in]{geometry}
\usepackage{booktabs,longtable,array}
\usepackage{microtype}
\usepackage[hidelinks]{hyperref}
\hypersetup{pdftitle={Report 154 Diagonal poly Cauchy permutations at all fixed orders},pdfauthor={},pdfsubject={OEIS A192563 exact saddle expansion, growing cycle moments, and inverse thresholds}}
\setlength{\parskip}{4pt}
\setlength{\parindent}{1.2em}
\setlength{\emergencystretch}{2em}
\setcounter{tocdepth}{1}
\numberwithin{equation}{section}
\newtheorem{theorem}{Theorem}[section]
\newtheorem{lemma}[theorem]{Lemma}
\newtheorem{proposition}[theorem]{Proposition}
\newtheorem{corollary}[theorem]{Corollary}
\theoremstyle{definition}\newtheorem{definition}[theorem]{Definition}
\theoremstyle{remark}\newtheorem{remark}[theorem]{Remark}
\newcommand{\EE}{\mathbb E}
\newcommand{\NN}{\mathbb N}
\newcommand{\QQ}{\mathbb Q}
\newcommand{\eps}{\varepsilon}
\newcommand{\sone}[2]{\genfrac{[}{]}{0pt}{}{#1}{#2}}
\newcommand{\stwo}[2]{\genfrac{\{} {\}}{0pt}{}{#1}{#2}}
\newcommand{\ceil}[1]{\left\lceil #1\right\rceil}
\newcommand{\floor}[1]{\left\lfloor #1\right\rfloor}
\newcommand{\dd}{\,\mathrm d}
\newcommand{\ii}{\mathrm i}
\title{Diagonal poly Cauchy permutations at all fixed orders\\\large Exact saddle expansion and inverse thresholds for OEIS A192563}
\author{Report 154}
\date{3 October 2026}
\begin{document}
\maketitle
\begin{abstract}
The diagonal second-kind poly-Cauchy numbers satisfy
$a_n=\sum_m\sone nm(m+1)^n$ and count the $(n,n)$-poly-Cauchy
permutations of B\'enyi and Ram\'irez. Equivalently,
$a_n/n!$ is the $n$th ordinary moment of one plus the cycle count
of a uniform permutation of $n$ labels. We prove a complete
asymptotic expansion at every fixed order, using the exact saddle
of $e^z\Gamma(n+e^z)/\Gamma(e^z)$ and explicit finite Gaussian
contraction polynomials in its standardized cumulants. The small scale is
$r/n$, where
$r=\log n-\log\log n-\log\log\log n+o(1)$.
A positive exponential sum controls the whole complementary
contour, and a weighted Taylor argument proves the integrated
remainder. The result is uniform when the moment order $k$ and
permutation size $n$ have ratio in a fixed compact subset of
$(0,\infty)$. Differentiating an explicit real-parameter model
then gives arbitrarily fine smooth inverse windows, with
strictly justified floor and ceiling bounds at equality cases.
We also obtain an explicit nested-logarithm inversion,
including its relative remainder. Exact companion calculations
check the coefficient identities and the OEIS prefix; optional
high-precision diagnostics are separated from the proof.
\end{abstract}
\noindent\textbf{Scope and attribution.}
The sequence, signed conventions, exact identities, and combinatorial
interpretation are established prior art. This report supplies a
sequence-specific fixed-order saddle analysis and its inverse
consequences. It does not claim exhaustive priority. Several related
full texts, notably Yakymiv's cycle-moment paper, remain uninspected.
No numerical asymptotic onset or universal error constant is certified
here. No claim concerns growing truncation order or exponentially
small sectors.
\newpage
\tableofcontents
\newpage

\section{Results and mathematical purpose}\label{sec:results}

Let $\sone nm$ and $\stwo nm$ be, respectively, the unsigned
Stirling numbers of the first kind and the Stirling numbers of the
second kind, with both $(0,0)$ entries equal to one. Define
\begin{equation}\label{eq:definition}
 A_{n,k}=\sum_{m=0}^n\sone nm(m+1)^k,
 \qquad a_n=A_{n,n}\quad(n,k\geq0).
\end{equation}
The diagonal begins
\[
1,\ 2,\ 13,\ 161,\ 3148,\ 87784,\ 3274640,\ 156359874.
\]
It is OEIS A192563, the diagonal of A344639
\cite{oeisdiag,oeisarray}. Throughout, logarithms are natural and
all asymptotic statements refer to positive arguments tending to
infinity.

\subsection{Why this diagonal is a natural asymptotic problem}

If $C_n$ is the number of cycles of a uniform random permutation
of $n$ labels, then
\begin{equation}\label{eq:moment}
 \frac{A_{n,k}}{n!}=\EE(C_n+1)^k.
\end{equation}
Thus the diagonal studies a moment order that grows with the
underlying combinatorial size. The familiar fixed-order cycle-moment
regime does not settle it. Large moments give disproportionate
weight to atypically large cycle counts, so replacing $C_n$ by its
typical order $\log n$ is not a justified approximation.

There is also a direct combinatorial class. Take a permutation of
$n$ red labels, put its cycles into canonical order with the largest
entry of each cycle first and these largest entries increasing,
and assign each of $k$ blue labels to one of the $m+1$ gaps around
its $m$ cycles. Within each gap put the blue labels in increasing
order. Removing cycle parentheses produces the second-kind
poly-Cauchy permutations of B\'enyi--Ram\'irez
\cite[Theorem 2.1]{benyi}. This description explains both positivity
and the factor $(m+1)^k$. It counts structured permutations on two
specified label sets, not arbitrary permutations of their total size.

The present results consequently serve three concrete purposes:
asymptotic enumeration of that class, evaluation of growing ordinary
cycle-count moments, and inversion of a prescribed counting or moment
scale. The compact aspect-ratio theorem permits the two label sets
to grow at comparable but unequal rates.

\subsection{The exact saddle and its coefficients}

For real $N\geq1$ and real $z$, set
\begin{equation}\label{eq:LN}
 L_N(z)=z+\log\Gamma(N+e^z)-\log\Gamma(e^z).
\end{equation}
Let $r=r(N)>0$ solve
\begin{equation}\label{eq:saddle}
 rL_N'(r)=N,
 \quad\text{or}\quad
 r\bigl[1+e^r\{\psi(N+e^r)-\psi(e^r)\}\bigr]=N.
\end{equation}
Here $\psi$ is the digamma function. Existence and uniqueness for
every $N\geq1$ are proved in Section~\ref{sec:real}. Define
\begin{equation}\label{eq:cumulants}
 \kappa_j(N)=\left.(z\partial_z)^jL_N(z)\right|_{z=r(N)},
 \qquad b=\kappa_2,
 \qquad h=\frac rN.
\end{equation}
The differential operator acts with $N$ held fixed, before evaluation
at the saddle. In particular $\kappa_1=N$ and $b>0$.

Put $E_0=1$. For $q\geq1$, define the finite polynomial expression
\begin{equation}\label{eq:Ek}
 E_q=\sum_{\substack{m_3,\ldots,m_{2q+2}\geq0\\
                    \sum_{j=3}^{2q+2}(j-2)m_j=2q}}
 (-1)^{J/2}(J-1)!!
 \prod_{j=3}^{2q+2}\frac1{m_j!}
       \left(\frac{\kappa_j}{j!b^{j/2}}\right)^{m_j},
 \qquad J=\sum_{j=3}^{2q+2}jm_j.
\end{equation}
The constraint makes $J$ even; the apparent half-powers of $b$ thus
combine to integer powers in every monomial. We use $(-1)!!=1$
when representing $E_0$ by the empty product.

\begin{theorem}[All fixed orders on the diagonal]\label{thm:main}
For every fixed integer $R\geq1$, as $n\to\infty$,
\begin{equation}\label{eq:main}
 a_n=\frac{n!\exp L_n(r)}{r^n\sqrt{2\pi b}}
       \left(\sum_{q=0}^{R-1}E_q(n)+O_R\left((r/n)^R\right)\right).
\end{equation}
Moreover, for each fixed $j\geq1$,
\begin{align}
 r(N)&=\log N-\log\log N-\log\log\log N+o(1),\label{eq:rscale}\\
 \kappa_j(N)&\sim Nr^{j-1},\qquad
 E_q(N)=O_q\left((r/N)^q\right).\label{eq:kappascale}
\end{align}
In particular,
\begin{equation}\label{eq:E1}
 E_1=\frac{\kappa_4}{8b^2}
       -\frac{5\kappa_3^2}{24b^3}
       \sim-\frac{r}{12N}.
\end{equation}
The implicit constants and sufficiently large onsets may depend on
$R$ or on the displayed fixed index.
\end{theorem}

The leading equivalent already involves an exact transcendental
saddle. Retaining it is essential at higher orders. The replacement
$E_1\sim-r/(12n)$ explains a sign and a scale, but does \emph{not}
replace the full $E_1$ in the $R=2$ formula. The errors in that
replacement, and in a finite nested-log approximation to $r$, can
be much larger than $(r/n)^2$. This is an exact-saddle expansion with
slowly varying coefficients, not a constant-coefficient power series
in $1/n$.

\begin{theorem}[Compact aspect ratios]\label{thm:aspect}
Fix $0<\alpha\leq\beta<\infty$ and a fixed integer $R\geq1$.
For positive integers $n,k$ with $\alpha\leq k/n\leq\beta$,
let $r>0$ solve $rL_n'(r)=k$ and define $\kappa_j,b,E_q$ as in
\eqref{eq:cumulants}--\eqref{eq:Ek}, with $n$ held fixed. Uniformly
in this range,
\begin{equation}\label{eq:aspect}
 A_{n,k}=\frac{k!\exp L_n(r)}{r^k\sqrt{2\pi b}}
    \left(\sum_{q=0}^{R-1}E_q
          +O_{R,\alpha,\beta}\left((r/k)^R\right)\right).
\end{equation}
Also, uniformly in the same range,
\begin{align}
 r&=\log n-\log\log n-\log\log\log n
               +\log(k/n)+o(1),\label{eq:aspectr}\\
 \kappa_j&\sim kr^{j-1},\qquad b\sim kr.
\end{align}
The theorem makes no uniform assertion as $k/n$ approaches zero or
infinity.
\end{theorem}

\subsection{Inverse results and rounding}

For fixed $R\geq1$, put
\begin{equation}\label{eq:YR}
 P_R(N)=\sum_{q=0}^{R-1}E_q(N),\qquad
 Y_R(N)=\frac{\Gamma(N+1)\exp L_N(r(N))}
                   {r(N)^N\sqrt{2\pi b(N)}}P_R(N).
\end{equation}
This defines an explicit smooth function for sufficiently large real
$N$. It does not assert an extension of the combinatorial sequence
to noninteger indices.

\begin{theorem}[Smooth inverse and integer thresholds]\label{thm:inverse}
The function $P_R$ is eventually positive and $Y_R$ is eventually
strictly increasing, with
\begin{equation}\label{eq:Yderivative}
 (\log Y_R)'(N)=2\log N-\log r(N)+o(1)\sim2\log N.
\end{equation}
For sufficiently large $X$, let $\nu_R$ be the unique solution of
$Y_R(\nu_R)=X$ in the eventual increasing range. There is a constant
$C_R>0$ such that, with
\begin{equation}\label{eq:epsilon}
 \eps_R=C_R\frac{(\log\nu_R)^{R-1}}{\nu_R^R},
\end{equation}
the two thresholds
\[
 T_{\geq}(X)=\min\{n\geq0:a_n\geq X\},\qquad
 T_{>}(X)=\min\{n\geq0:a_n>X\}
\]
satisfy
\begin{equation}\label{eq:bracket}
 \floor{\nu_R-\eps_R}+1
 \ \leq T_{\geq}(X)\leq T_{>}(X)
 \ \leq\ceil{\nu_R+\eps_R}.
\end{equation}
In addition, if $L=\log X$, $\ell=\log L$, and $m=\log\ell$,
then
\begin{equation}\label{eq:nested}
 \nu_R=\frac{L}{2\ell}
       \left[1+\frac{\tfrac32m+1+\log2}{\ell}
               +O_R\left(\frac{m^2}{\ell^2}\right)\right].
\end{equation}
\end{theorem}

The shrinking $\eps_R$ measures uncertainty around a \emph{smooth}
center before rounding. An integer threshold cannot in general equal
$\nu_R+o(1)$ without a rounding qualification. If the interval
$[\nu_R-\eps_R,\nu_R+\eps_R]$ lies strictly between consecutive
integers $j$ and $j+1$, both thresholds equal $j+1$. If $X=a_n$,
then $T_{\geq}(X)=n$ and $T_>(X)=n+1$; the two-sided statement
\eqref{eq:bracket} deliberately retains that possibility.
The explicit nested-log expression has a much coarser error than
\eqref{eq:epsilon}; it must not silently replace $\nu_R$ in that
fine bracket.

\section{Established identities and source boundaries}\label{sec:sources}

\subsection{Signs and the OEIS identification}

In Komatsu's notation for second-kind poly-Cauchy numbers,
\[
 A_{n,k}=(-1)^n\widehat c_n^{(-k)}.
\]
B\'enyi--Ram\'irez write the nonnegative normalization as
$\widehat c_{n,k}$. These conventions exclude the signed first-kind
poly-Cauchy numbers, which are a different family
\cite[Theorem 4]{komatsu2013}. The final published version of
B\'enyi--Ram\'irez gives \eqref{eq:definition} and its interpretation
in Theorem 2.1 and gives
\begin{equation}\label{eq:stirlingproduct}
 A_{n,k}=\sum_{j=0}^k j!\sone{n+1}{j+1}\stwo{k+1}{j+1}
\end{equation}
in Theorem 2.6 \cite{benyi}. The diagonal of
\eqref{eq:stirlingproduct} is precisely the defining formula in
A192563. The theorem numbers in the earlier preprint are different;
all theorem references here use the final 2022 volume version.

For completeness, a direct algebraic check of
\eqref{eq:stirlingproduct} is short. The standard identity
\[
 (m+1)^k=\sum_{j=0}^k\stwo{k+1}{j+1}(m)_j
\]
uses the falling factorial $(m)_j=m(m-1)\cdots(m-j+1)$.
The coefficient of $x^j$ in
\[
 \sum_{m=0}^n\sone nm(1+x)^m
      =(1+x)(2+x)\cdots(n+x)
\]
is $\sone{n+1}{j+1}$, and also equals
$\sum_m\sone nm\binom mj$. Multiplying by $j!$ and interchanging
finite sums proves the formula. The companion verifies both sides
by independent recurrences.

\subsection{The coefficient function}

The rising-factorial identity gives the entire function
\begin{align}
 F_n(z)&=\sum_{m=0}^n\sone nm e^{(m+1)z}\label{eq:positive}\\
       &=e^z\prod_{q=0}^{n-1}(e^z+q)
        =e^z\frac{\Gamma(n+e^z)}{\Gamma(e^z)}.
        \label{eq:Fn}
\end{align}
The gamma quotient in \eqref{eq:Fn} is understood through its
finite product wherever necessary; $F_n$ is entire. Hence
\begin{equation}\label{eq:coefficient}
 A_{n,k}=k![z^k]F_n(z),\qquad a_n=n![z^n]F_n(z).
\end{equation}
Summing the rising-factorial exponential generating function in
$n$ also gives
\begin{equation}\label{eq:doubleegf}
 \sum_{n,k\geq0}A_{n,k}\frac{x^n}{n!}\frac{y^k}{k!}
     =e^y(1-x)^{-e^y}.
\end{equation}
These are exact prior-art identities, including the shifted-family
specialization behind \eqref{eq:doubleegf}. The one-variable form
\eqref{eq:coefficient} is the convenient route for the asymptotic
analysis; it changes with $n$, so the proof must control a triangular
family of entire functions.

\subsection{What the literature screen establishes}

The official OEIS data-mirror entry for A192563 was inspected in
full on 3 October 2026, as was A344639; the latter was also read at
its OEIS page. Neither inspected entry contains a diagonal asymptotic.
The individual A192563 webpage and its b-file were not successfully
retrieved in that pass. Thus the verified finite prefix used here
comes from the official mirror, not from a claimed b-file inspection.

Full inspected primary texts include Komatsu's 2013 article
\cite{komatsu2013}, the final B\'enyi--Ram\'irez article
\cite{benyi}, Komatsu's 2021 transform preprint
\cite{transform}, and Jha's cycle-moment preprint \cite{jha}.
The first three provide exact identities and related generalizations;
no matching diagonal all-orders theorem was located in those texts.
Jha's fixed-order factorial-moment result is not a theorem uniform in
the growing ordinary moment order used here.

The following limits matter to any novelty claim. The full texts of
Komatsu--Szalay's shifted-family paper \cite{shifted}, Komatsu's
2019 recurrence paper \cite{recurrence}, and the final 2022 version
of the transform paper were not inspected. Most importantly,
Yakymiv's paper on moments of cycle counts of restricted permutations
\cite{yakymiv} was identified through its official abstract, but its
full text and its moment-order uniformity were not verified. This
report therefore does not rule out overlap with that paper. A
bounded search for the sequence identifiers, diagonal poly-Cauchy
asymptotics, and growing cycle-count moments found no matching
result; absence from that search is not an exhaustive priority proof.

The analytic inputs below are classical gamma-function formulas
\cite{dlmf}. The mathematical contribution presented here is their
application with a complete fixed-order error argument, exact
coefficient retention, compact aspect-ratio uniformity, and
rounding-aware inverse consequences. Source snapshots and precise
inspection levels are recorded in the accompanying provenance file.

\section{Real parameter saddles and scale estimates}\label{sec:real}

This section works with real $N\geq1$. It simultaneously supplies
integer saddle estimates and the smooth continuation required for
inversion. Write $t=e^r$ and
\begin{equation}\label{eq:q}
 q_N(t)=t\{\psi(N+t)-\psi(t)\}.
\end{equation}

\begin{lemma}[Existence and uniqueness]\label{lem:exist}
For every real $N\geq1$, the equation $rL_N'(r)=N$ has exactly
one solution $r>0$. At this solution $b>0$, and $r(N)$ is smooth.
\end{lemma}
\begin{proof}
The digamma difference integral from DLMF 5.9.16 \cite{dlmf} gives
\begin{equation}\label{eq:digammadiff}
 q_N(t)=\int_0^\infty te^{-tu}h_N(u)\dd u,
 \qquad h_N(u)=\frac{1-e^{-Nu}}{1-e^{-u}}.
\end{equation}
For $u>0$,
\[
 \frac{h_N'(u)}{h_N(u)}
 =\frac{N}{e^{Nu}-1}-\frac1{e^u-1}\leq0,
\]
because $e^{Nu}-1\geq N(e^u-1)$ for $N\geq1$.
The last inequality follows, for example, from convexity in the
exponent parameter. If $V$ is exponential with rate one,
\eqref{eq:digammadiff} is $\EE h_N(V/t)$. Since $h_N$ is
nonincreasing, this expectation is nondecreasing in $t$.
The integral is smooth for $t>0$, so $q_N'(t)\geq0$.
Consequently
\[
 L_N'(r)=1+q_N(e^r)>0,\qquad L_N''(r)\geq0.
\]
Thus $rL_N'(r)$ is strictly increasing for $r>0$, begins at zero,
and tends to infinity, since it is at least $r$. Its derivative is
$L_N'(r)+rL_N''(r)=b/r>0$. Existence, uniqueness, and smoothness
follow from continuity and the implicit function theorem.
\end{proof}

\begin{lemma}[Derivative bounds with an integral remainder]\label{lem:psi}
For every fixed integer $j\geq0$ and $x>0$,
\begin{equation}\label{eq:psiremainder}
 \left|\frac{\mathrm d^j}{\mathrm dx^j}
          \{\psi(x)-\log x\}\right|
       \leq \frac{j!}{x^{j+1}}.
\end{equation}
In particular, for each fixed $j\geq1$,
\begin{equation}\label{eq:polygamma}
 \psi(x)=\log x+O(x^{-1}),\qquad
 \psi^{(j)}(x)=(-1)^{j-1}(j-1)!x^{-j}+O_j(x^{-j-1}).
\end{equation}
\end{lemma}
\begin{proof}
The digamma remainder integral, DLMF 5.9.13 \cite{dlmf}, is
\begin{equation}\label{eq:G}
 \psi(x)-\log x=-\int_0^\infty G(u)e^{-xu}\dd u,
 \qquad G(u)=\frac1{1-e^{-u}}-\frac1u.
\end{equation}
For $u>0$ one has $0<G(u)<1$: the lower inequality uses
$1-e^{-u}<u$, and the upper one uses $e^u-1>u$.
On each compact subinterval of $x>0$, dominated differentiation is
valid to every fixed order. The absolute value of the $j$th
derivative of \eqref{eq:G} is at most
$\int_0^\infty u^je^{-xu}\dd u=j!x^{-j-1}$.
This proves the claims without differentiating an unspecified
asymptotic remainder.
\end{proof}

\begin{lemma}[Saddle and cumulant scales]\label{lem:scales}
As $N\to\infty$, with $t=e^{r(N)}$,
\begin{equation}\label{eq:tscale}
 \frac tN\sim\frac1{r\log r},\qquad
 r\sim\log N,\qquad
 r=\log N-\log\log N-\log\log\log N+o(1).
\end{equation}
For every fixed $j\geq1$,
\begin{equation}\label{eq:Lderivatives}
 L_N^{(j)}(r)=\frac Nr+O_j(t+1)\sim\frac Nr,
 \qquad \kappa_j\sim Nr^{j-1}.
\end{equation}
\end{lemma}
\begin{proof}
Lemma~\ref{lem:psi} first yields, uniformly for $N,t\geq1$,
\begin{equation}\label{eq:qapprox}
 q_N(t)=t\log(1+N/t)+O(1).
\end{equation}
At $r=\tfrac12\log N$, the left side of the saddle equation is
$O(\sqrt N\log^2N)=o(N)$. At $r=\log N$ it exceeds $N$ for
large $N$, since $q_N(N)=N\log2+O(1)$. Monotonicity therefore
gives
\begin{equation}\label{eq:roughr}
 \tfrac12\log N<r<\log N
\end{equation}
for all sufficiently large $N$.

Put $x=N/t$. The saddle equation and \eqref{eq:qapprox} give
\begin{equation}\label{eq:xequation}
 \frac xr=\log(1+x)+O(t^{-1}).
\end{equation}
Here $x\geq1$, $r\to\infty$, and $t\to\infty$.
A bounded subsequence of $x$ would contradict
\eqref{eq:xequation}, so $x\to\infty$. Hence
$x/(r\log x)\to1$. The function $x/\log x$ is eventually
increasing. Substitution of fixed multiples of $r\log r$ into
this function brackets $x$ between such multiples, proving
$\log x\sim\log r$ and then $x\sim r\log r$.
Taking logarithms once more gives
\[
 \log x=\log r+\log\log r+o(1).
\]
Since $r=\log N-\log x$, the refined formula in
\eqref{eq:tscale} follows from \eqref{eq:roughr} and
$r/\log N\to1$.

For derivatives, apply $(t\partial_t)^{j-1}$ to $q_N(t)$.
The explicit bounds \eqref{eq:psiremainder} imply that the
remainder in \eqref{eq:qapprox} and each of its fixed
$t\partial_t$ derivatives are $O_j(1)$. Indeed each resulting
term has a power of $t$ multiplying a derivative bounded at
$t$ or at $N+t$ by the corresponding inverse power. In the main
term $t\log(1+N/t)$, the original logarithmic term survives
once after every application of $t\partial_t$; all other terms
are $O_j(t)$. It follows that
\[
 L_N^{(j)}(r)=t\log(1+N/t)+O_j(t+1)
           =N/r+O_j(t+1).
\]
For $j=1$ the saddle equation gives the exact equality
$L_N'(r)=N/r$. Since $rt/N\sim1/\log r\to0$, these expressions
are asymptotic to $N/r$.

Finally, the Stirling differential-operator identity is
\begin{equation}\label{eq:operator}
 (z\partial_z)^jL_N(z)
   =\sum_{v=1}^j\stwo jv z^v L_N^{(v)}(z).
\end{equation}
At $z=r$, its term $v=j$ is asymptotic to $Nr^{j-1}$ and its
other fixed terms are smaller. This proves the cumulant estimate.
\end{proof}

\subsection{A finite Bernoulli representation at integer size}

For integer $n$, the product form gives a second, elementary
calculus for the derivatives. Write
\begin{equation}\label{eq:p}
 p_q=\frac{t}{t+q}\quad(0\leq q\leq n-1).
\end{equation}
Let $B_j(p)$ be the $j$th cumulant polynomial of a Bernoulli
variable with success probability $p$. These polynomials satisfy
\begin{equation}\label{eq:bernoulli}
 B_1(p)=p,\qquad
 B_{j+1}(p)=p(1-p)B_j'(p),\qquad B_j(p)=p+O_j(p^2).
\end{equation}
Then
\begin{equation}\label{eq:integerderivatives}
 L_n'=1+\sum_qp_q,\qquad
 L_n^{(j)}=\sum_qB_j(p_q)\quad(j\geq2).
\end{equation}
Integral comparison shows
\[
 \sum_qp_q=t\log(1+n/t)+O(1),\qquad
 \sum_qp_q^2=O(t+1).
\]
This recovers the fixed-derivative estimates without numerical
differentiation. It also furnishes the optional diagnostic program
with a numerically stable finite-sum route, independent of repeated
differentiation of the gamma quotient.

\section{Contour control and the integrated expansion}\label{sec:contour}

Only integer $n$ is needed in the coefficient integral. All real
parameter arguments above serve estimates and the later inverse;
we do not integrate a noninteger-size generating function.
Fix a small constant $0<\rho<\pi/3$.

\begin{lemma}[A fixed complex neighborhood]\label{lem:complex}
For each fixed $j\geq1$, the analytic logarithm of $F_n$ agreeing
with $L_n$ on the real line satisfies
\begin{equation}\label{eq:complex}
 |L_n^{(j)}(z)|\leq C_jn/r\qquad(|z-r|\leq\rho)
\end{equation}
for all sufficiently large $n$. If
$f(\theta)=L_n(re^{\ii\theta})$, then, on intervals with
$r|\theta|=o(1)$,
\begin{equation}\label{eq:fderivatives}
 |f^{(j)}(\theta)|\leq C_jnr^{j-1}.
\end{equation}
\end{lemma}
\begin{proof}
In the disk $|z-r|\leq\rho$, one has
$\operatorname{Re}e^z\geq c e^r=ct$. Hence every factor
$e^z+q$ is nonzero there, and
\[
 \left|\frac{e^z}{e^z+q}\right|
       \leq C\frac{t}{t+q}=Cp_q.
\]
The logarithmic derivatives of $e^z+q$ are the same fixed
Bernoulli cumulant polynomials as in \eqref{eq:bernoulli},
evaluated at $e^z/(e^z+q)$. Each has zero constant coefficient.
The sum of their absolute values is bounded by a fixed multiple
of $\sum_qp_q=O(n/r)$; the extra term $z$ is harmless.
A nonvanishing analytic function on this disk has an analytic
logarithm, establishing \eqref{eq:complex}. Applying the chain
rule repeatedly, or \eqref{eq:operator}, gives
\eqref{eq:fderivatives}.
\end{proof}

Set $\delta=n^{-2/5}$. At $\theta=0$,
$f^{(j)}(0)=\ii^j\kappa_j$. Taylor's theorem through degree two
therefore gives, uniformly for $|\theta|\leq\delta$,
\begin{align}
 \operatorname{Re}\{L_n(re^{\ii\theta})-L_n(r)-\ii n\theta\}
 &\leq-\tfrac12b\theta^2+Cnr^2|\theta|^3\notag\\
 &\leq-\tfrac14b\theta^2.\label{eq:localgaussian}
\end{align}
The last inequality follows from $b\sim nr$ and $r\delta\to0$.

\begin{lemma}[The full complementary circle]\label{lem:minor}
For some fixed $c>0$ and all large $n$,
\begin{equation}\label{eq:minor}
 \sup_{\delta\leq|\theta|\leq\pi}
 \frac{|F_n(re^{\ii\theta})|}{F_n(r)}
       \leq e^{-cn\delta^2}=e^{-cn^{1/5}}.
\end{equation}
\end{lemma}
\begin{proof}
The positive exponential sum \eqref{eq:positive} implies
\begin{align*}
 |F_n(re^{\ii\theta})|
 &\leq\sum_m\sone nm e^{(m+1)r\cos\theta}\\
 &=F_n(r\cos\theta)\leq F_n(r\cos\delta)
       \qquad(\delta\leq|\theta|\leq\pi).
\end{align*}
The function $F_n$ is increasing on the entire real line because
all its exponential-sum coefficients are nonnegative and their
exponents are positive. On $s\in[r\cos\delta,r]$, one has
$s-r=o(1)$. The finite-sum estimates
\eqref{eq:integerderivatives} give $L_n'(s)\sim n/r$ uniformly
there. Since $r-r\cos\delta\asymp r\delta^2$,
\[
 L_n(r\cos\delta)-L_n(r)
       =-\int_{r\cos\delta}^r L_n'(s)\dd s
       \leq-cn\delta^2.
\]
This proves the assertion over the whole complementary circle.
\end{proof}

The last argument is useful because $e^z$ has an imaginary period,
which can make an informal local saddle calculation look vulnerable
to secondary peaks. The bound handles the entire circle at once;
no alleged periodic peak has been omitted.

\subsection{Normalization and the central cutoff}

Cauchy's coefficient formula is
\begin{equation}\label{eq:cauchy}
 \frac{a_nr^n}{n!F_n(r)}
 =\frac1{2\pi}\int_{-\pi}^{\pi}
      \exp\{L_n(re^{\ii\theta})-L_n(r)-\ii n\theta\}\dd\theta,
\end{equation}
where the integrand on the full circle is interpreted as
$F_n(re^{\ii\theta})e^{-\ii n\theta}/F_n(r)$, so a global
logarithm is unnecessary. Only the local arc uses its chosen
analytic logarithm. After multiplying by $\sqrt{2\pi b}$,
Lemma~\ref{lem:minor} bounds the complementary contribution by
\[
 O\bigl(\sqrt b\,e^{-cn^{1/5}}\bigr),
\]
which is smaller than every fixed power of $r/n$.

Write
\begin{equation}\label{eq:normalization}
 \eta=\sqrt{r/n},\qquad u=\sqrt b\,\theta,\qquad
 \lambda_j=\frac{\kappa_j}{b^{j/2}}.
\end{equation}
Then $\lambda_j=O_j(\eta^{j-2})$ for every fixed $j\geq3$.
Restrict further to $|u|\leq U=\eta^{-1/6}$. This interval lies
inside $|\theta|\leq\delta$, since
\[
 \frac U{\sqrt b}=O(n^{-5/12}r^{-7/12})=o(n^{-2/5}).
\]
By \eqref{eq:localgaussian}, the omitted part of the local arc
has normalized integral $O(e^{-cU^2})$. This too is smaller than
any fixed power of $\eta$.

\subsection{A weighted remainder that can be integrated}

Fix $R\geq1$. Taylor expansion through degree $2R+2$, using
Lemma~\ref{lem:complex} at degree $2R+3$, gives
\begin{equation}\label{eq:centralphase}
 L_n(re^{\ii\theta})-L_n(r)-\ii n\theta
       =-u^2/2+Q(u)+\mathcal R(u),
\end{equation}
where
\begin{equation}\label{eq:Q}
 Q(u)=\sum_{j=3}^{2R+2}\frac{\lambda_j(\ii u)^j}{j!},
 \qquad
 |\mathcal R(u)|\leq C_R\eta^{2R+1}|u|^{2R+3}.
\end{equation}
Indeed the normalized derivative bound is
$nr^{2R+2}/b^{(2R+3)/2}=O_R(\eta^{2R+1})$.
On $|u|\leq U$, the bounds imply
$|Q(u)|\leq C_R\eta|u|^3=o(1)$ and
$|\mathcal R(u)|=o(1)$. Therefore
\[
 |e^{Q(u)+\mathcal R(u)}-e^{Q(u)}|
       \leq C_R|\mathcal R(u)|,
\]
and its integral against $e^{-u^2/2}$ is
$O_R(\eta^{2R+1})$.

For the remaining exponential, write
$\lambda_j=\eta^{j-2}c_j$, where all the finitely many $c_j$
are bounded uniformly in $n$. Introduce an auxiliary real variable
$s$ and set
\begin{equation}\label{eq:H}
 H(s,u)=\exp\left(
    \sum_{j=3}^{2R+2}\frac{c_j(\ii u)^j s^{j-2}}{j!}
                       \right).
\end{equation}
The exponent is bounded uniformly for $0\leq s\leq\eta$ and
$|u|\leq U$. Differentiating $H$ exactly $2R$ times with respect
to $s$ produces $H$ times a finite polynomial in $u$, the bounded
$c_j$, and nonnegative powers of $s$. Thus there is a fixed
integer $M_R$ for which Taylor's theorem in $s$ yields
\begin{equation}\label{eq:weightedTaylor}
 H(\eta,u)=\sum_{d=0}^{2R-1}\eta^dH_d(u)
          +O_R\bigl(\eta^{2R}(1+|u|^{M_R})\bigr).
\end{equation}
The bound is integrable against the Gaussian, uniformly in $n$.
This is the necessary integrated error argument. A pointwise Taylor
bound evaluated only at the outer cutoff would not by itself prove
the asserted relative remainder.

Every monomial of weighted degree $d$ has ordinary degree
$J=\sum j m_j$ in $u$, while
$d=\sum(j-2)m_j$. Hence $J\equiv d\pmod2$.
Odd $d$ terms integrate to zero on the symmetric interval.
For even $d=2q$, extending each polynomial Gaussian integral from
$[-U,U]$ to $\mathbb R$ incurs only an exponentially decaying
error. The identities
\begin{equation}\label{eq:Gaussian}
 \frac1{\sqrt{2\pi}}\int_{\mathbb R}e^{-u^2/2}u^{2v}\dd u
     =(2v-1)!!,\qquad
 \int_{\mathbb R}e^{-u^2/2}u^{2v+1}\dd u=0
\end{equation}
give precisely the contraction formula \eqref{eq:Ek}.
Combining \eqref{eq:cauchy}--\eqref{eq:weightedTaylor} proves
\eqref{eq:main} with error $O_R(\eta^{2R})=O_R((r/n)^R)$.
Each monomial of $E_q$ is $O_q(\eta^{2q})$, proving its size
claim. Finally, substituting $\kappa_j\sim nr^{j-1}$ and
$b\sim nr$ into \eqref{eq:E1} gives
$(1/8-5/24)r/n=-r/(12n)$. This completes the proof of
Theorem~\ref{thm:main}.

\section{Exact coefficient calculus}\label{sec:coefficients}

The coefficients are finite universal expressions in the exact
cumulants. For example, the next one after \eqref{eq:E1} is
\begin{align}\label{eq:E2}
 E_2={}&-\frac{\kappa_6}{48b^3}
      +\frac{7\kappa_3\kappa_5}{48b^4}
      +\frac{35\kappa_4^2}{384b^4}\notag\\
     &-\frac{35\kappa_3^2\kappa_4}{64b^5}
      +\frac{385\kappa_3^4}{1152b^6}.
\end{align}
No limiting replacement of a cumulant is part of this formula.

\subsection{Two independent finite generation rules}

Formula \eqref{eq:Ek} can be generated by enumerating the
nonnegative integer solutions of
$\sum_{j\geq3}(j-2)m_j=2q$. There is also an algebraically
independent formal exponential recurrence. For indeterminates
$x_3,x_4,\ldots$, set
\begin{equation}\label{eq:formalQ}
 Q_w(u)=\frac{x_{w+2}(\ii u)^{w+2}}{(w+2)!},\qquad
 \exp\left(\sum_{w\geq1}Q_w(u)s^w\right)
     =\sum_{d\geq0}H_d(u)s^d.
\end{equation}
Formal differentiation with respect to $s$ gives
\begin{equation}\label{eq:formalrecurrence}
 H_0(u)=1,\qquad
 H_d(u)=\frac1d\sum_{w=1}^d wQ_w(u)H_{d-w}(u).
\end{equation}
Integrating $H_{2q}$ term by term against the standard Gaussian
and then putting $x_j=\kappa_j/b^{j/2}$ yields $E_q$.
Every step takes place in a finite rational polynomial ring after
a degree cutoff; convergence of the formal infinite series is not
asserted or needed. The companion compares this recurrence with
the weighted-partition formula using exact rational arithmetic.

For the Bernoulli polynomials there are likewise two exact rules:
\eqref{eq:bernoulli} and
\begin{equation}\label{eq:BernoulliStirling}
 B_j(p)=\sum_{v=1}^j(-1)^{v-1}(v-1)!\stwo jv p^v.
\end{equation}
The latter follows by expanding
$\log(1+p(e^z-1))$ in powers of $p$ and extracting the $j$th
derivative at zero. Comparison of these two constructions checks
the derivative input used by the numerical diagnostics.

\subsection{What an exact finite check certifies}

The rational contraction identities, Stirling recurrences,
Bernoulli identities, and agreement of finite sequence values can
all be certified by terminating exact calculations. They check
signs, factorial normalizations, parity, and the correct sequence
convention. They do not prove the uniform contour estimate or an
asymptotic remainder. Those claims depend on the analytic arguments
in Sections~\ref{sec:real} and \ref{sec:contour}.

Similarly, a small numerical residual at a finite list of sizes is
not an interval bound for all larger sizes. The optional numerical
program is explicitly a diagnostic. In particular, it cannot turn
the asymptotic threshold theorem into a certified finite-$X$
rounding oracle without explicit constants and an onset.

\section{Uniformity for comparable moment orders}\label{sec:aspect}

We give the details needed for Theorem~\ref{thm:aspect}, rather than
merely substituting $k$ for $n$ in the diagonal formula. Let
$\lambda=k/n\in[\alpha,\beta]$ and choose the unique positive
saddle satisfying $rL_n'(r)=k$. The same existence proof applies,
since $rL_n'(r)$ increases from zero to infinity.

The preliminary comparison at $r=\tfrac12\log n$ and
$r=\log n$ now gives
\[
 \tfrac12\log n<r<\log n
\]
for all large $n$, uniformly in $\lambda$. Put $t=e^r$ and
$x=n/t$. The saddle relation becomes
\begin{equation}\label{eq:aspectx}
 \frac{\lambda x}{r}=\log(1+x)+O(t^{-1}).
\end{equation}
Since $\lambda$ stays bounded above and away from zero,
$x\to\infty$ uniformly and inversion gives
\[
 x\sim\frac r\lambda\log r,
 \qquad \log x=\log r+\log\log r-\log\lambda+o(1),
\]
uniformly. Taking $r=\log n-\log x$ proves
\eqref{eq:aspectr}. The derivative estimates become
\[
 L_n^{(j)}(r)=k/r+O_j(t+1)\sim k/r,
 \qquad \kappa_j\sim kr^{j-1}.
\]
The constants can be chosen uniformly on the compact ratio interval.

The disk $|z-r|\leq\rho$ used in Lemma~\ref{lem:complex} is
unchanged. The derivative bound is now $C_jkr^{j-1}$ for the
angular function. With $\delta=n^{-2/5}$, positivity gives a
complementary-circle bound $\exp(-ck\delta^2)$, where $c>0$
can be chosen uniformly in $\lambda$. Since $k\asymp n$, this
is at most $\exp(-c' n^{1/5})$. The central Gaussian argument
uses $\eta=\sqrt{r/k}$ and $U=\eta^{-1/6}$; both the cutoff
inclusion and every fixed-degree integrated remainder remain
uniform. Applying Cauchy's formula to
$A_{n,k}=k![z^k]F_n(z)$ proves \eqref{eq:aspect}.

This theorem leaves two different boundary regimes open. Fixed
or slowly growing $k$ is not covered by a compact positive ratio;
neither is $k/n\to\infty$. A theorem connecting these regimes
would need new uniform estimates, not an unqualified appeal to
\eqref{eq:aspect}.

\section{Differentiating the smooth model}\label{sec:differentiation}

We now prove the derivative claim needed for inversion. All
derivatives in this section are of exact gamma, polygamma, and
finite cumulant expressions. No derivative is taken of the
asymptotic error in \eqref{eq:main}.

\begin{lemma}[Parameter derivatives]\label{lem:parameter}
For fixed $j\geq1$, along the real saddle,
\begin{equation}\label{eq:kappaderivative}
 r_N\sim N^{-1},\qquad
 \frac{\mathrm d\kappa_j}{\mathrm dN}=O_j(r^{j-1}),
 \qquad \frac{b_N}{b}=O(N^{-1}).
\end{equation}
For fixed $q\geq1$,
\begin{equation}\label{eq:Ederivative}
 E_q'(N)=O_q\left(\frac{(r/N)^q}{N}\right).
\end{equation}
\end{lemma}
\begin{proof}
With $r$ temporarily held fixed, direct differentiation gives
\begin{equation}\label{eq:partialN}
 \partial_NL_N^{(j)}(r)
  =(t\partial_t)^j\psi(N+t)
  =\sum_{v=1}^j\stwo jv t^v\psi^{(v)}(N+t)
  =O_j(t/N).
\end{equation}
The last bound follows from Lemma~\ref{lem:psi}, since $t/N\to0$.
Implicit differentiation of \eqref{eq:saddle} gives exactly
\begin{equation}\label{eq:rprime}
 r_N=\frac{r\{1-rt\psi'(N+t)\}}{b}.
\end{equation}
Now $rt\psi'(N+t)\sim rt/N\sim1/\log r\to0$ and
$b\sim Nr$, proving $r_N\sim1/N$.

From \eqref{eq:operator}, \eqref{eq:Lderivatives}, and
\eqref{eq:partialN}, with the other variable fixed,
\[
 \partial_r\kappa_j=O_j(Nr^{j-1}),\qquad
 \partial_N\kappa_j=O_j(r^jt/N).
\]
The chain rule and \eqref{eq:rprime} give
\[
 \frac{\mathrm d\kappa_j}{\mathrm dN}
 =\partial_N\kappa_j+r_N\partial_r\kappa_j
 =O_j(r^{j-1}),
\]
because $rt/N=O(1)$. Division by
$\kappa_j\sim Nr^{j-1}$ shows that every fixed $\kappa_j$
is eventually positive with logarithmic derivative $O_j(1/N)$.
The same applies to $b$.

Differentiate each monomial of \eqref{eq:Ek} separately.
Its size is $O_q((r/N)^q)$, and the product rule multiplies
that bound by $O_q(1/N)$. Summing finitely many terms proves
\eqref{eq:Ederivative}. This argument does not assume that the
sum $E_q$ itself is nonzero.
\end{proof}

It follows that
\begin{equation}\label{eq:Pbounds}
 P_R=1+O_R(r/N),\qquad
 P_R'=O_R(r/N^2).
\end{equation}
For $R=1$, the derivative is exactly zero. Thus $P_R>0$ for
large $N$. In the logarithm of \eqref{eq:YR}, all terms involving
$r_N$ from $L_N(r)-N\log r$ cancel by the saddle equation.
The exact derivative is
\begin{equation}\label{eq:envelope}
 (\log Y_R)'=
 \psi(N+1)+\psi(N+t)-\log r
 -\frac{b_N}{2b}+\frac{P_R'}{P_R}.
\end{equation}
Lemma~\ref{lem:psi} and \eqref{eq:Pbounds} yield
\begin{equation}\label{eq:logYderivative}
 (\log Y_R)'=2\log N-\log r
             +O_R(t/N+1/N)
           =2\log N-\log r+o(1).
\end{equation}
This proves eventual strict monotonicity. It also shows
$Y_R(N)\to\infty$, so every sufficiently large $X$ has one
solution in the eventual increasing range. There is no need to
assert global monotonicity of $Y_R$ at small $N$.

\section{Threshold windows and strict equality edges}\label{sec:thresholds}

We first check exact monotonicity of the original sequence.
The unsigned Stirling recurrence gives, for $n\geq1$,
\begin{equation}\label{eq:monotone}
 a_{n+1}=\sum_{m=0}^n\sone nm
       \{n(m+1)^{n+1}+(m+2)^{n+1}\}>a_n.
\end{equation}
Also $a_1=2>a_0=1$. Since $a_n\geq n!$, the sequence tends to
infinity, and both thresholds are well defined.

Theorem~\ref{thm:main} and $P_R=1+O_R(r/n)$ imply
\begin{equation}\label{eq:logerror}
 \log a_n-\log Y_R(n)
   =O_R\left((r(n)/n)^R\right)
   =O_R\left((\log n/n)^R\right).
\end{equation}
Fix a sufficiently large $X$, abbreviate $\nu=\nu_R$, and take
\begin{equation}\label{eq:rounded}
 n_- =\floor{\nu-\eps},\qquad
 n_+ =\ceil{\nu+\eps},\qquad
 \eps=C_R\frac{(\log\nu)^{R-1}}{\nu^R}.
\end{equation}
Both integers are asymptotic to $\nu$ and, for large $X$, lie
beyond every onset used in the proof. On the short intervals
joining them to $\nu$, \eqref{eq:logYderivative} is at least
$\log\nu$. Consequently
\begin{align}
 \log Y_R(\nu)-\log Y_R(n_-)&\geq\eps\log\nu,\notag\\
 \log Y_R(n_+)-\log Y_R(\nu)&\geq\eps\log\nu.
 \label{eq:strictgaps}
\end{align}
The local errors in \eqref{eq:logerror} are bounded by a constant
multiple of $(\log\nu/\nu)^R$. Choose $C_R$ strictly larger
than that constant, increasing the onset if necessary. Since
$\eps\log\nu=C_R(\log\nu/\nu)^R$, this gives the
\emph{strict} inequalities
\begin{equation}\label{eq:strictendpoints}
 a_{n_-}<X<a_{n_+}.
\end{equation}
Exact monotonicity now proves \eqref{eq:bracket}. A fixed initial
prefix causes no difficulty: increasing $X$ past its largest entry
excludes it automatically.

\subsection{Why the strict inequalities are necessary}

The lower edge is $n_-+1$, not merely $n_-$. The upper edge
bounds the strict threshold because $X<a_{n_+}$, not just
$X\leq a_{n_+}$. If $X=a_j$, the threshold definitions give
$j$ and $j+1$ respectively. Any derivation which tacitly identifies
these two thresholds at equality would be false. The strictly
larger error constant in \eqref{eq:strictgaps} ensures that
\eqref{eq:bracket} covers this case as written.

The theorem supplies a hierarchy of centers. At $R=1$, the
smooth uncertainty is $O(\nu^{-1})$; at $R=2$ it is
$O(\log\nu/\nu^2)$; at $R=3$ it is
$O((\log\nu)^2/\nu^3)$. The constants are not numerically
certified here. These are asymptotic precision statements, not
precomputed finite-input guarantees. A certified implementation
would require effective constants for the contour and central
errors together with interval evaluation of the saddle and inverse.

\section{First corrected nested logarithm inversion}\label{sec:nested}

The fine inverse center is defined implicitly, but its leading
nested-logarithm form can be proved with an explicit relative
remainder. We begin with the required logarithmic growth estimate.

\begin{lemma}[Logarithmic size of the smooth model]\label{lem:loggrowth}
For every fixed $R\geq1$,
\begin{equation}\label{eq:loggrowth}
 \log Y_R(N)
 =2N\log N-N\log\log N-2N
   +O_R\left(\frac{N\log\log N}{\log N}\right).
\end{equation}
\end{lemma}
\begin{proof}
The logarithmic Stirling formula with its positive-real remainder
bound \cite[5.11.1 and 5.11(ii)]{dlmf}, together with $t/N\to0$,
gives
\begin{equation}\label{eq:Stirlingratio}
 \log\Gamma(N+t)-\log\Gamma(t)
 =N\log N-N+t\log(N/t)
      +O(t+t^2/N+\log N).
\end{equation}
For example, expand $(N+t)\log(N+t)$ around $N$ using
$\log(1+t/N)=t/N+O((t/N)^2)$; the remaining half-logarithm
terms are $O(\log N)$ and the Stirling remainders are bounded.
The saddle relation yields
$t\log(1+N/t)=N/r+O(1)$ and hence
$t\log(N/t)=O(N/r)$. Also
$t+t^2/N+\log N=O(N/r)$.
Adding the initial term $r$ in \eqref{eq:LN} therefore gives
\[
 L_N(r)=N\log N-N+O(N/r).
\]
Now use
$\log\Gamma(N+1)=N\log N-N+O(\log N)$,
$\log b=O(\log N)$, and $\log P_R=O_R(r/N)$ in
\eqref{eq:YR}. We obtain
\begin{equation}\label{eq:loggrowthr}
 \log Y_R(N)=2N\log N-N\log r-2N+O_R(N/r).
\end{equation}
From \eqref{eq:rscale},
$\log r=\log\log N+O(\log\log N/\log N)$.
Substitution proves \eqref{eq:loggrowth}.
\end{proof}

Let $L=\log X$, $\ell=\log L$, and $m=\log\ell$.
Equation \eqref{eq:loggrowth} first implies
$\nu_R\sim L/(2\ell)$: it gives
$L\sim2\nu_R\log\nu_R$, so
$\log\nu_R\sim\log L$ and then the stated equivalent.
Set
\begin{equation}\label{eq:x0}
 c=\tfrac32m+1+\log2,
 \qquad x_0=\frac{L}{2\ell}\left(1+\frac c\ell\right).
\end{equation}
Elementary logarithmic expansions give
\begin{align}
 \log x_0&=\ell-m-\log2+c/\ell+O(m^2/\ell^2),
                  \label{eq:logx0}\\
 \log\log x_0&=m-(m+\log2)/\ell+O(m^2/\ell^2).
                  \label{eq:loglogx0}
\end{align}
On substituting \eqref{eq:logx0}--\eqref{eq:loglogx0} into
\eqref{eq:loggrowth}, the coefficient of $L/\ell$ cancels:
\[
 c-\tfrac32m-1-\log2=0.
\]
The result is
\begin{equation}\label{eq:x0error}
 \log Y_R(x_0)=L+O_R(Lm^2/\ell^2).
\end{equation}
Both $x_0$ and $\nu_R$ are asymptotic to $L/(2\ell)$.
Throughout the interval between them, \eqref{eq:Yderivative}
is of order $\ell$. The mean value theorem applied to
\eqref{eq:x0error} consequently gives
\[
 \nu_R-x_0=O_R(Lm^2/\ell^3).
\]
This is exactly \eqref{eq:nested}, including its relative
$O_R(m^2/\ell^2)$ remainder, and completes the proof of
Theorem~\ref{thm:inverse}.

The coefficient $3/2$ has two contributions: inversion of
$N\log N$ contributes one factor of $\log\log L$, and the
term $-N\log\log N$ contributes an additional half after
normalization by $2N\log N$. The constants $1+\log2$ likewise
come from the $-2N$ term and the leading factor two. This explanation
is an interpretation of the explicit cancellation, not a substitute
for its remainder calculation.

\section{Reproducible exact checks and numerical diagnostics}\label{sec:repro}

The accompanying archive contains this PDF, its standalone LaTeX
source, deterministic build and packaging scripts, source provenance,
and a companion directory. It does not require a private proof file,
an earlier report, or a live network connection to verify its exact
identities. Only the optional diagnostics require an external Python
package, pinned as \texttt{mpmath==1.3.0}.

The exact companion checks are designed around independent routes:
\begin{enumerate}
\item Generate unsigned first-kind and second-kind Stirling rows by
      their integer recurrences, and compare the moment sum with
      \eqref{eq:stirlingproduct} and a derivative-product evaluation
      of \eqref{eq:Fn}.
\item Reproduce the entire 17-term prefix in the inspected official
      A192563 mirror snapshot, with offset zero and the specified
      sign normalization.
\item Compare weighted-partition Gaussian contractions with the
      formal exponential recurrence \eqref{eq:formalrecurrence},
      using exact rational coefficients. Check the displayed
      $E_1$ and $E_2$ separately.
\item Compare Bernoulli cumulant polynomials generated by
      \eqref{eq:bernoulli} and \eqref{eq:BernoulliStirling}, and
      verify the Stirling differential-operator identity at finite
      degrees.
\item Check strict and nonstrict threshold definitions at exact
      finite equality cases, and exercise the output-safety rules
      under both ordinary Python and \texttt{python3 -O}.
\end{enumerate}
The released exact run checks all $4225$ pairs $0\leq n,k\leq64$
by the three integer routes and constructs $E_0,\ldots,E_6$.
Their monomial counts are $1,2,5,11,22,42,77$. It performs
$11993$ individual finite checks, including the polynomial and
threshold identities; normal and optimized runs have the same
mathematical payload. No correctness condition relies on Python's
removable \texttt{assert} statements. The companion's documented finite
cutoffs are verification parameters, not hypotheses of the theorems.
The companion README lists the commands; result files record the checks.

Optional diagnostics solve the real saddle, evaluate cumulants by
finite Bernoulli sums, and compare the leading and corrected saddle
models with exact integer counts. They report the normalized error
\begin{equation}\label{eq:diagnostic}
 \frac{a_n/M_n-\sum_{q=0}^{R-1}E_q}{(r/n)^R},
 \qquad M_n=\frac{n!F_n(r)}{r^n\sqrt{2\pi b}},
\end{equation}
and analogous compact-aspect-ratio samples. These values are
floating-point observations, separated from exact pass/fail
identities and from the proof of the remainder. The diagnostics do
not infer rigorous interval bounds or an asymptotic onset from
sampled residuals.


\begin{table}[htbp]
\centering
\caption{Observed scaled residuals in \eqref{eq:diagnostic}.
These rounded 100-digit calculations are diagnostics, not error bounds.}
\label{tab:residuals}
\begin{tabular}{rrrr}
\toprule
$n$ & $R=2$ & $R=3$ & $R=4$\\
\midrule
20 & -0.008764515 & -0.001372810 & 0.000515071 \\
50 & -0.006294384 & 0.000461294 & 0.001454634 \\
100 & -0.004675041 & 0.001130859 & 0.001471817 \\
200 & -0.003449208 & 0.001464946 & 0.001341716 \\
500 & -0.002300129 & 0.001667662 & 0.001143029 \\
1000 & -0.001672183 & 0.001739809 & 0.001010160 \\
\bottomrule
\end{tabular}
\end{table}

The numerical file additionally covers aspect ratios $1/2$ and $2$
at $n=50,200,500$, and all corrections through $E_6$. Changes of
sign in the scaled residuals are compatible with the theorem:
it supplies a fixed-order magnitude estimate, not a prescribed sign
or a monotone approach to a limit.

\subsection{Build and release boundaries}

The scripts use exclusive new output paths and reject static symlinked
path components. Builds assume a trusted parent directory with no
concurrent malicious path replacement; directory creation is not a
security boundary against a same-user process swapping the new leaf
before it is opened. PDF builds use a fresh directory, fixed metadata,
no shell escape, and stable cross-references. The release packager
reads a fixed allowlist, checks every member against the SHA-256
manifest, and creates a stored ZIP with fixed timestamps, ordering,
and permissions. Byte reproducibility of the PDF is checked in
clean builds under the recorded TeX toolchain; a different TeX or
font version may produce different PDF bytes. Exact mathematical
checks need only the Python standard library.

No publication, OEIS edit, repository modification, or external
submission is part of this report. The source references are public;
the archive distributes original exposition and programs, not copies
of the cited journal PDFs or private review material.

\section{Limits and further research}\label{sec:questions}

The results cover each fixed truncation order with a proved
integrated remainder. They do not provide estimates uniform in
$R\to\infty$, an optimal truncation theorem, or a decomposition
into exponentially small asymptotic sectors. The fact that the
noncentral contour is smaller than every fixed algebraic order
does not, by itself, supply any of those stronger claims.

Several extensions have concrete mathematical content:
\begin{enumerate}
\item \textbf{Effective remainder constants.}
      Can one turn the disk, contour, and weighted Taylor bounds
      into practical explicit constants and onsets? Combining
      those with interval gamma and polygamma evaluation could
      produce certified finite-$X$ threshold decisions.
\item \textbf{Transition between moment regimes.}
      Determine a uniform expansion as $k/n\to0$, connecting
      fixed-order cycle moments to the present growing-order
      regime. The transition as $k/n\to\infty$ is a separate
      problem; it may involve substantially different saddle scales.
\item \textbf{Cycle statistics under the moment weight.}
      For the probability proportional to
      $\sone nm(m+1)^k$, obtain a limit law and local estimates
      for $m$. This would describe the cycle counts typical of
      large poly-Cauchy structures rather than just their total
      number. No such limit law is asserted here.
\item \textbf{Further explicit inverse terms.}
      Sharpen \eqref{eq:loggrowthr} using higher saddle-scale and
      gamma expansions, and derive a full hierarchy in nested
      logarithms with proved remainders. Its error should remain
      distinct from the much finer exact-saddle inverse window.
\item \textbf{Restricted and shifted families.}
      Investigate restricted cycle lengths and shifted poly-Cauchy
      parameters after a complete comparison with Yakymiv and the
      shifted-family literature. Positivity may still control
      global contours, but the local derivative scales require
      separate proofs.
\item \textbf{Large truncation orders.}
      Study the growth of the contraction polynomials and their
      analytic remainder constants. A valid growing-order or
      exponential-improvement theorem would need information beyond
      the fixed-order estimates established here.
\end{enumerate}

These questions are proposed research directions, not claims of
proved results or of absence from all prior literature. In particular,
reading the uninspected cycle-moment source is a priority before
making stronger originality statements about restricted permutations
or growing moments.

\begin{thebibliography}{99}

\bibitem{oeisdiag}
OEIS Foundation Inc., \emph{A192563}, official OEIS data mirror,
entry last changed 26 May 2021, inspected 3 October 2026.
\href{https://github.com/oeis/oeisdata/blob/main/seq/A192/A192563.seq}
{Official mirror entry}.
The inspected snapshot digest is recorded in the source provenance file.

\bibitem{oeisarray}
OEIS Foundation Inc., \emph{A344639}, sign-normalized second-kind
poly-Cauchy array, inspected 3 October 2026.
\href{https://oeis.org/A344639}{OEIS entry};
\href{https://github.com/oeis/oeisdata/blob/main/seq/A344/A344639.seq}
{official mirror}.

\bibitem{benyi}
B. B\'enyi and J. L. Ram\'irez,
\emph{Poly-Cauchy Numbers of the Second Kind---the Combinatorics Behind},
Enumerative Combinatorics and Applications \textbf{2}:1 (2022),
Article S2R1, 13 pp.
\href{https://doi.org/10.54550/ECA2022V2S1R1}{DOI};
\href{https://ecajournal.haifa.ac.il/Volume2022/ECA2022_S2A1.pdf}
{inspected final PDF}.
Final Theorems 2.1 and 2.6 are used here. The publisher records
online publication in June 2021 for the 2022 volume.

\bibitem{komatsu2013}
T. Komatsu, \emph{Poly-Cauchy numbers},
Kyushu Journal of Mathematics \textbf{67} (2013), 143--153.
\href{https://doi.org/10.2206/kyushujm.67.143}{DOI};
\href{https://www.jstage.jst.go.jp/article/kyushujm/67/1/67_143/_pdf}
{inspected final PDF}.

\bibitem{transform}
T. Komatsu, \emph{Recurrence relations of poly-Cauchy numbers by
the $r$-Stirling transform}, arXiv:2103.15291 (2021),
subsequently Mediterranean Journal of Mathematics \textbf{19} (2022).
\href{https://arxiv.org/abs/2103.15291}{Inspected preprint}.
A separate full final version was not inspected.

\bibitem{jha}
S. K. Jha, \emph{Asymptotic expansions for factorial moments of
some distributions in the analysis of algorithms},
arXiv:1611.07336 (2016).
\href{https://arxiv.org/abs/1611.07336}{Inspected preprint}.
The cycle-count moment order in Theorem 2.2 is fixed.

\bibitem{shifted}
T. Komatsu and L. Szalay, \emph{Shifted poly-Cauchy numbers},
Lithuanian Mathematical Journal \textbf{54}:2 (2014), 166--181.
\href{https://doi.org/10.1007/s10986-014-9235-y}{DOI}.
Only repository and indexed extracts were inspected, including the
bivariate formula associated with Proposition 3; full text uninspected.

\bibitem{recurrence}
T. Komatsu, \emph{Some recurrence relations of poly-Cauchy numbers},
Journal of Nonlinear Sciences and Applications \textbf{12}:12
(2019), 829--845.
\href{https://www.isr-publications.com/jnsa/articles-8077-some-recurrence-relations-of-poly-cauchy-numbers}
{Publisher record}. Abstract and indexed theorem only; full text
uninspected.

\bibitem{yakymiv}
A. L. Yakymiv, \emph{Asymptotics with remainder term for moments
of the total cycle number of random A-permutation},
Diskretnaya Matematika \textbf{31}:3 (2019), 114--127;
English translation, Discrete Mathematics and Applications
\textbf{31}:1 (2021), 51--60.
\href{https://doi.org/10.4213/dm1582}{Russian DOI};
\href{https://doi.org/10.1515/dma-2021-0005}{translation DOI}.
Official abstract inspected; full text and uniformity in moment order
unverified.

\bibitem{dlmf}
NIST Digital Library of Mathematical Functions, Chapter 5,
\emph{Gamma Function}, equations
\href{https://dlmf.nist.gov/5.9.E13}{5.9.13},
\href{https://dlmf.nist.gov/5.9.E16}{5.9.16},
\href{https://dlmf.nist.gov/5.11.E1}{5.11.1}, and
\href{https://dlmf.nist.gov/5.11.ii}{Section 5.11(ii)};
verified 3 October 2026. The fixed-derivative bound used here is
derived directly from the integral remainder. Standard fuller
expansions also appear in
\href{https://dlmf.nist.gov/5.11.E2}{5.11.2} and
\href{https://dlmf.nist.gov/5.15.E9}{5.15.9}.

\end{thebibliography}
\end{document}
